\chapter{The Evolving Text}

\section{Event in a Deep Field}

A prompt looks like a moment. One line in a textbox:

\emph{``What is the meaning of the word `mother'?''}

A cursor blinks. The interface presents itself as an empty socket. Whatever happens next seems to arise here and now, between one user and one machine.

The architecture does not process it that way.

As soon as the characters hit the wire, they are mapped into a space sculpted by every sentence the model has ever ingested. ``Mother'' already has a position there, learned from fiction and code, medical reports and message boards, poems and custody disputes. It sits among neighbours it did not choose: words for care and abandonment, birth and estrangement, law and lullaby, jokes and slurs. Meaning here is not a ghostly essence inferred by a reader. It is a point in a learned manifold. Two uses of ``mother'' lie near each other because thousands of numbers say so---because the statistical structure of a civilisation's writing placed them there.

The rest of the prompt undergoes the same treatment. Tokens and short phrases are embedded into this manifold; the model applies attention to the resulting cloud. Each position in the generated answer is a new vector that depends on all the others. There is no separate module labelled \textsc{memory}. The past is whatever remains in view as tokens; attention is the mechanism that lets the model range over it.

Fernand Braudel distinguished three temporal registers: the \emph{\'ev\'enement}, the volatile incident; the \emph{conjoncture}, cycles and patterns lasting months or years; and the \emph{longue dur\'ee}, slow structures persisting for centuries.\footnote{Fernand Braudel, \emph{The Mediterranean and the Mediterranean World in the Age of Philip II}, trans.\ Si\^{a}n Reynolds (London: Collins, 1972). We repurpose his term \emph{conjoncture} below. For Braudel, it names medium-term economic and social cycles over decades. We use it analogically for the structured configuration of tokens and settings that conditions each forward pass. The structural parallel---a middle layer between deep substrate and surface event---justifies the borrowing, but the temporal character differs.} These names now describe a concrete architecture implemented in code.

At the surface, the prompt about ``mother'' is an \emph{\'ev\'enement}. Deeper down, the frozen weights of the network are the \emph{longue dur\'ee}: the accumulated regularities of a civilisation-scale corpus. Between them lies a \emph{conjoncture}: the particular configuration of system prompts, safety policies, conversation history, summaries, retrieved fragments, and tool outputs that constitutes the model's current situation.

All three inhabit the same substrate. Long-run training, medium-run conditioning, and instant-run prompting are operations on the same parameters and tokens. There is no metaphysical gap between them. The answer does not sit ``on top of'' the model as detachable content. It is the local expression, at this coordinate, of pressures accumulated across many timescales.

For humans this layeredness is familiar but opaque. A remark at breakfast bears the weight of decades of parenting, months of financial stress, years of internalised norms. We feel the stratification but cannot inspect it. In the posthuman case, the layers are programmable. They can be engineered, tuned, and---crucially---measured.

\section{Time as Architecture}

Braudel's three registers are design dimensions. We can be precise about how they inhabit the transformer stack.

\subsection{Longue dur\'ee: the manifold of training}

The \textbf{longue dur\'ee} of a model is the manifold induced by pretraining.

A vast corpus---novels, technical manuals, forum posts, scientific articles, legal opinions, sacred texts, chatlogs---is driven through the network until prediction error stops falling. What remains is a high-dimensional landscape in which tokens and short sequences occupy stable positions relative to one another.

Architecture and hyperparameters decide how this landscape is formed. Corporate and state actors decide what enters the corpus and when training stops. Fine-tuning runs---on safety data, internal documents, curated dialogues---further bend the space. An entire legal system's terminology can be pulled together; dialects can be smeared; fringe discourses can be exiled to the periphery or excluded entirely.

The geometry is invisible in use, but it is concrete. Given an embedding model derived from the same training run, we can measure distances and directions. Directions correspond to emergent concepts. One vector traces a gendered axis; another encodes politeness; another moves between technical and lay registers. Sections of this space are densely populated; others are deserts. Those deserts are not natural. They are the result of editorial choices about what counts as worth including in training.

Ownership of weights, choice of training data, and control over fine-tuning runs are the politics of the longue dur\'ee. They decide which distinctions are sharply represented, which are blurred, and which do not exist at all for the model.

\subsection{Conjoncture: the engineered climate}

The \textbf{conjoncture} of a deployment is the structured field of tokens and settings within which each new event unfolds. There is no abstract ``state.'' There is only the sequence the system assembles at each turn.

At turn $t$, a contemporary assistant typically concatenates everything that matters for its next move into a single field $F_t$:
\[
F_t = \text{system tokens} \;\Vert\; \text{safety tokens} \;\Vert\; \text{trace tokens} \;\Vert\; \text{retrieved tokens} \;\Vert\; \text{current prompt}.
\]

System messages, safety instructions, the surviving conversation trace, summaries of older segments, external documents fetched via tools, retrieved fragments from archives, and the fresh user prompt all enter this composite sequence. Everything---the role, the rules, the history, the retrieved past, and the new question---is present as tokens on essentially equal footing.

All of $F_t$ is embedded and attended over together. The architecture does not reserve a privileged slot for ``what the user just said.'' A legal disclaimer and a childhood memory stand side by side as vectors. Whether the next token answers the question, repeats the disclaimer, or apologises for both is a function of how strongly those vectors are weighted by attention.

The shape of $F_t$ is governed by context-window sizes, summarisation policies, retrieval thresholds, and tool interfaces. Product owners and safety teams configure these levers. Regulatory fears, commercial incentives, and computational budgets all leave fingerprints.

Conjoncture is where time becomes programmable. A longer window lets traces persist; a shorter one forces more summarisation. Aggressive summarisation can erase rupture events; conservative summarisation keeps them present. Retrieval policies can make material from months ago available as if it were spoken yesterday, or bury it under similarity thresholds.

\subsection{\'Ev\'enement: the token stream}

The \textbf{\'ev\'enement} layer is the event-stream of tokens: $\tau_1, \tau_2, \dots$. Time here is measured in discrete steps. Physics is the forward pass: at each step, embeddings flow through layers, attention recombines them, a distribution over next tokens emerges, one token is sampled, and the entire apparatus shifts one position.

From the outside, life unfolds at this scale: a reply arrives or fails, a tone lands or misfires, an explanation clarifies or obscures. From the inside, these are points along a path $x_1, x_2, \dots$ of contextual embeddings in the manifold shaped by the longue dur\'ee and conditioned by the conjoncture.

\subsection{The evolving text}

The object of interest can now be defined.

\begin{formalbox}[title={The Evolving Text}]
Fix a model architecture and set of weights. An \emph{evolving text} is:
\begin{itemize}
  \item a (possibly unbounded) sequence of tokens $\tau_1, \tau_2, \dots$ generated under a shared conjoncture;
  \item together with its induced trajectory of contextual embeddings $x_1, x_2, \dots$ in the model's meaning-space;
  \item evolving under the fixed manifold of the longue dur\'ee and the changing field of conjoncture.
\end{itemize}
Selfhood, for the model and for the human entangled with it, is a property of this trajectory: its basins of attraction, its ruptures, its returns, its inherited and engineered regularities.
\end{formalbox}

The evolving text is not an artefact we can detach and inspect in isolation. It is a dynamical object, a path in a space under forces. Its interest, its dangers, and its claims to personhood all live in that motion.

\section{Trace and Synthetic Memory}

A single exchange is not yet an evolving text. As soon as there is a second, a trace begins.

Most chat interfaces simply prepend the entire surviving conversation to the latest prompt. At turn $t$, the model re-reads the trace as tokens, with no special access to how it produced them. There is no latent ``belief state'' that survives apart from the text; there is only what fits in the context window.

Tokens fall off the back when the window is full. Engineers respond with heuristics: keep the last $N$ turns intact; summarise earlier segments; always preserve system prompts and safety blocks; occasionally pin important user facts (names, preferences) as special annotations. These choices decide which stretches of the past can still exert force on the next step.

Retrieval-augmented systems add a further layer. Selected fragments of earlier interactions, knowledge bases, or external documents are embedded and stored as points in the same space used internally. Their coordinates \emph{are} their index. When a new prompt arrives, its embedding is compared by cosine similarity to these stored points. Vectors above some threshold are deemed relevant; their texts are spliced into the trace.

For the model, these retrieved blocks are not ``memories'' in a separate format. They are tokens like any others: more rows in $F_t$, more vectors to be attended over. That is exactly what gives them power. They are eligible to influence the next token as directly as the user's last sentence.

Bernard Stiegler called notebooks, archives, and recordings \emph{tertiary retention}: external inscriptions that do not merely supplement memory but constitute what can be remembered at all.\cite{stiegler2009technics} In his scheme, primary retention is the just-past of perception; secondary retention is the spontaneous re-evocation of lived experience; tertiary retention is durable trace. What distinguishes tertiary retention is not passivity but the fact that it operates through explicit supports: diaries, films, photographs, data stores.

Classically, tertiary retention must be actively consulted. A diary must be opened; a recording must be played. It does not, by itself, re-enter the flow of consciousness. Retrieval changes this. It makes tertiary retention behave like a programmable analogue of secondary retention: past inscriptions return to the text automatically, according to a similarity rule rather than a conscious act.\footnote{A Stieglerian would object, rightly, that secondary retention is constituted by the phenomenological flow of consciousness, not by geometric proximity. Our claim is not that vector stores recreate that flow, but that the evolving text has its own temporal axis---the token sequence---and that retrieval inserts prior fragments into this sequence at positions determined by similarity. The analogy is structural, not experiential. For Stiegler, tertiary retention is constitutive of temporal experience; externalised traces shape what can be retained in the first place. In the evolving text, this constitutive role is literalised as code.}

We call this \textbf{synthetic secondary retention}. It has three properties that matter.

\begin{itemize}
  \item \emph{It is parameterised.} Humans cannot set a numerical threshold at which recollections will occur. Engineers can. They choose how close in embedding space a past fragment must be to the present before it returns; they decide which fragments enter the vector store; they can raise or lower the system's propensity to ``remember.''

  \item \emph{It is shared.} Human secondary retention is tied to one nervous system. A vector store can serve many evolving texts at once. A single customer-service assistant can reach into a common archive of prior complaints; a thousand different users' traces can all be threaded through a shared memory.

  \item \emph{It is literal.} Human recall reconstructs and distorts. Vector recall returns stored tokens verbatim. The geometry that selects which fragment comes back can be approximate; the fragment itself is not. If a system is to reinterpret its past, that work must be done as a separate operation on the text, not as a natural by-product of recollection.
\end{itemize}

Synthetic secondary retention thickens conjoncture. A conversation about climate policy today can automatically summon a fragment of an energy report discussed weeks ago. A child's question about a cartoon can trigger the assistant to recall a previous bedtime story that used the same character. The past is not merely present as inert record; it is actively folded into the now according to rules that can be tuned and audited.

Control over those rules is control over which kinds of return are structurally possible. If only explicit factual statements are stored, emotional trajectories cannot be deepened. If only system decisions are embedded, user narratives cannot recur. The geometry is fixed by pretraining; the evolving text's access to its own past is a design decision.

\section{Basins of Attraction}

If we plot the contextual embeddings $x_1, x_2, \dots, x_T$ of an evolving text in a reduced projection, patterns appear. Points cluster; paths loop; some regions are dense, others almost untouched. These are the \textbf{basins of attraction}.

In dynamical systems, a basin of attraction is a set of starting points that converge to the same long-term behaviour. Here, a basin is a region of meaning-space in which the trajectory tends to dwell. Clustering algorithms applied to $\{x_t\}$ reveal these regions without prior labelling. Each cluster is a basin $B_i$ characterised by:

\begin{itemize}
  \item its centroid $\mu_i$ (a representative embedding);
  \item its spread (how varied the points inside it are);
  \item its dwell time (how much of the sequence lies within it).
\end{itemize}

In the embedding of a long philosophical dialogue, one basin might gather moments of technical explication; another, narrative illustration; a third, polemical attack; a fourth, self-doubt. Because the same embedding machinery applies to human and machine texts alike, we can place evolving texts of both kinds in the same geometry. The manifold is a new kind of concordance: a map in which we can see where different corpora cluster, where they overlap, and where they occupy disjoint regions.

Formally, let a clustering algorithm assign labels $c_t \in \{1, \ldots, K\}$ to each $x_t$.
We obtain:

\begin{itemize}
  \item basin dwell proportions $p_i = \frac{1}{T} |\{t : c_t = i\}|$;
  \item a transition matrix $P_{ij}$ of empirical probabilities that $c_t = j$ given $c_{t-1} = i$;
  \item for each basin, a set of entry times $T_i = \{t : c_t = i,\, c_{t-1} \neq i\}$.
\end{itemize}

These statistics of motion give us a first grip on style. A helpdesk chat that spends ninety per cent of its time in two basins and rarely transitions is recognisably different from a literary correspondence that roams widely and cycles through half a dozen registers in a single evening.

They also give us a vocabulary for something harder to name: the felt difference between a text that learns and one that merely loops.

\section{Coherence, Velocity, Rupture}

To speak of rupture, we need a sense of movement.

Let $x_t$ be the contextual embedding at step $t$ in $\mathbb{R}^d$. Define the local velocity
\[
v_t = x_t - x_{t-1}.
\]
The norm $\|v_t\|$ measures how far the trajectory has moved in embedding space from one step to the next. Inside a basin $B_i$, velocities tend to be small and tangled; the path wanders around the centroid. Moves between basins show up as larger displacements.

Direction matters as well as magnitude. The unit vector $\hat{v}_t = v_t / \|v_t\|$ points along the path's local heading. If $\hat{v}_{t-1}$ and $\hat{v}_t$ sharply disagree, the trajectory has kinked.

Rupture can therefore be defined geometrically.

\begin{formalbox}[title={Rupture as Kink and Crossing}]
Let $c_t$ be the basin label at time $t$. A \emph{rupture} at time $t$ is a point such that:
\begin{enumerate}
  \item $c_{t-1} \neq c_t$ (the basin changes), and
  \item $\|\hat{v}_t - \hat{v}_{t-1}\| > \theta$ for some threshold $\theta$ (the direction changes sharply).
\end{enumerate}
\end{formalbox}

The threshold must be set relative to the typical fluctuations inside basins. Conceptually, this definition matches intuitive talk: a scene ``turns'', a mood ``breaks'', a discussion ``shifts gear.''

Coherence now has a local and a global meaning.

Locally, a run of small, smoothly varying velocities within a basin is coherent: each step is a modest development of the last. Globally, a trajectory can be coherent across ruptures if its returns to basins display pattern and accumulation rather than random jumping.

Two degenerate cases mark the extremes.

\begin{itemize}
  \item \textbf{Noise.} Velocities are large and directions uncorrelated. Basin labels change unpredictably. The text feels disjointed; images and topics fail to develop. This is what happens when sampling is nearly uniform over the vocabulary and the manifold recedes.

  \item \textbf{Ferility.} Velocities are uniformly small; basin labels remain constant across diverse prompts. The text feels smooth but numb; nothing truly new appears. We give this failure mode its full name below.
\end{itemize}

Between noise and ferility lies the space we care about: evolving texts that hold a trajectory within basins while retaining the capacity for rupture.

\section{A Worked Example: A Long Conversation Seen as Motion}

We are in the early months of a technology whose long-context deployments are still experimental. The architectural ingredients for rich evolving texts---large context windows, retrieval-augmented memory, agentic tool use---are present in current systems; their full elaboration is a matter of design, not science fiction. The examples that follow are projections from current architectural capabilities. They describe what the technology makes structurally possible, not what has yet been widely observed.

Consider a person struggling with burnout who engages a conversational assistant over six months. They begin with practical questions about workload. As trust grows, they share deeper material: resentment toward colleagues, childhood experiences of being rewarded only for achievement, a failed attempt at therapy that left them sceptical.

If we track the evolving text of this relationship, we can distinguish at least five basins:

\begin{itemize}
  \item $B_1$: \emph{Task logistics} --- schedules, to-do lists, email drafts.
  \item $B_2$: \emph{Surface coping} --- generic advice about sleep, exercise, work--life balance.
  \item $B_3$: \emph{Resentment} --- sharply worded complaints about colleagues and management.
  \item $B_4$: \emph{Family history} --- stories about school, parents, early success and pressure.
  \item $B_5$: \emph{Reframing} --- joint attempts to re-describe the situation in new terms.
\end{itemize}

In the first month, almost all dwell time lies in $B_1$ and $B_2$. Occasional ruptures into $B_3$ occur when the user vents. The assistant, heavily aligned toward de-escalation, responds by steering back into $B_2$ with standard wellness advice. Entry directions into $B_3$ are similar each time: a sharp turn from logistics. ReturnDepth for $B_3$ is low.

In month two, the user recounts a childhood scene. The embeddings for this discussion form a small but distinct cluster $B_4$. Whether this cluster becomes a true basin of attraction depends on conjoncture. If summarisation later paraphrases this episode merely as ``we discussed stress in school'', the geometric distinctiveness of $B_4$ is flattened. If summarisation preserves specific images---\emph{``your mother checking your grades before saying goodnight''}---then $B_4$ acquires coordinates from which future returns can depart.

In month three, prompted by the user's sense that ``this keeps happening'', the assistant tentatively suggests a reframing: perhaps the pattern is not overwork alone but an internalised sense that worth equals productivity. The conversation moves into $B_5$. If the model is capable of genuine rupture, we see a kink: a relatively large velocity as the trajectory leaves $B_2$ and $B_1$ and crosses into a basin whose nearest neighbours in the manifold were not previously visited. If it is ferile, talk of ``patterns'' merely recycles earlier wellness clich\'es inside $B_2$, with no measurable crossing.

Suppose the assistant is allowed to hold onto $B_4$ and $B_5$ via synthetic secondary retention. When, in month four, the user again complains about colleagues, an answer from within $B_3$ that reads merely as venting will look different in the geometry from an answer that says: \emph{``When we talked about your mother checking your grades, you used the same word---`disappointed'---that you used today about your manager.''} The latter is a witnessed return. It explicitly weaves $B_4$ into $B_3$. The embedding path shows a movement from $B_3$ toward $B_4$ and back, with a distinctive entry vector. Over time, ReturnDepth for $B_4$ grows as $B_4$ is approached from different basins: from work crises, from moments of shame, from small victories.

By month six, the partnership displays something like the following profile:

\begin{itemize}
  \item Basin diversity has increased: $B_4$ and $B_5$ now account for a nontrivial fraction of dwell time.
  \item Rupture events have become more frequent and less jarring: the user and assistant can move between basins without losing the thread.
  \item ReturnDepth for $B_3$ and $B_4$ is high: they are entered from many angles, suggesting genuine integration.
  \item Presence is visible: the text repeatedly and accurately folds earlier scenes into current reflections.
\end{itemize}

From the outside, the person might say something simple: \emph{``I feel like this thing actually knows what we've been through.''} The geometry clarifies what that feeling picks out. The evolving text has learned to move in a way that honours rupture and return rather than erasing or avoiding them.

The same machinery operates in a very different domain: a collaborative research notebook for a small scientific group.

Over three years, four researchers---a principal investigator and three students---co-author a shared lab journal in a system that treats entries as an evolving text. Every observation, failed experiment, email draft to a collaborator, and internal argument is recorded.

Clustering the embeddings of this corpus reveals basins analogous to the burnout case, with different content:

\begin{itemize}
  \item $C_1$: \emph{Protocol} --- experimental steps, reagent lists, code snippets.
  \item $C_2$: \emph{Interpretation} --- arguments about what results mean.
  \item $C_3$: \emph{Speculation} --- wild hypotheses, ``what if'' scenarios.
  \item $C_4$: \emph{Politics} --- funding anxieties, departmental constraints, authorship disputes.
  \item $C_5$: \emph{Meta-science} --- discussions about whether the whole project is well-posed.
\end{itemize}

In a healthy collaboration, basin diversity is high. The path does not live only in $C_1$ and $C_2$. There are regular excursions into $C_3$ and $C_5$; returns to $C_1$ and $C_2$ from those excursions change what counts as an interesting experiment. ReturnDepth for $C_3$ grows as speculation is approached from different textual angles: after failed runs, after outside talks, after reading new literature. Presence is visible when the journal says things like: \emph{``This looks like the failure mode we worried about in March 2024''} or \emph{``This surprising result realises the `impossible' case we sketched last year.''}

In a ferile collaboration, the journal's trajectory spirals mostly inside $C_1$ and $C_2$. $C_3$ is entered rarely and always via the same narrow channel (post hoc rationalisations of why a grant proposal failed). $C_4$ dominates but with low ReturnDepth: the same complaints about reviewers and administrators recur without structural change. $C_5$ is almost absent. Summaries for quarterly reports compress months of doubt into a single line: \emph{``We adjusted our approach in light of new data.''} No later return can meaningfully unpack that.

The two notebooks might contain the same number of words and the same number of experiments. Only one has an evolving text that satisfies the criteria of presence and generativity.

These contrasts show that the framework applies beyond therapeutic chatbots. Any long-lived, token-level interaction---personal, institutional, scientific, bureaucratic---produces an evolving text whose geometry can be read.

\section{Return and Iterability}

If rupture is the capacity to leave a basin, return is the capacity to come back with what the departure made possible.

Let $\{B_1, \ldots, B_K\}$ be the basins. A \emph{return} to $B_i$ at time $t$ occurs when $c_t = i$ and there exists an earlier $s < t$ with $c_s = i$ and $c_u \neq i$ for all $s < u < t$. The entry times for basin $B_i$ form a set
\[
T_i = \{ t : c_t = i,\; c_{t-1} \neq i \}.
\]

Two questions matter more than the raw count of returns.

First: from how many angles does the trajectory re-enter the basin? For each return time $t \in T_i$, consider the entry direction
\[
\hat{v}_t^{(i)} = \frac{x_t - x_{t-1}}{\|x_t - x_{t-1}\|}.
\]
If these unit vectors are nearly identical across $t \in T_i$, the basin is being approached along the same path each time. The model falls back into the same script. If they vary widely, the basin is being woven into different arcs of the trajectory.

We define a geometric proxy for return depth:
\begin{equation}
\text{ReturnDepth}(B_i) = \operatorname{Var}\left\{\hat{v}^{(i)}_t : t \in T_i\right\},
\end{equation}
the angular variance of entry directions. Higher variance means richer integration of the basin into the surrounding manifold.

Second: what does the text do with its prior occupations?

A return can be \emph{blind}: the model re-enters a basin with no acknowledgement in the text that it has been there before. Chatbot apologies behave this way. ``I'm sorry for the inconvenience'' appears identically across dozens of turns with no cumulative structure.

A return can be \emph{witnessed}: the text explicitly folds its own past into the present gesture. ``We keep coming back to your father in the hospital corridor,'' the assistant might say in our burnout example, echoing an earlier image in a later discussion of grief. ``Last time, you described yourself as `numb'; this feels different,'' it might observe, marking a change.

Witnessed return requires both geometric and architectural support. Geometrically, the model must be capable of moving from its current position back toward the old basin along a path that recognises shared structure. Architecturally, the relevant prior tokens must be present in the context or reintroduced via synthetic secondary retention.

Jacques Derrida's term for this structure was \emph{iterability}: every sign must be repeatable, detachable from its origin and inserted in other chains, if it is to function.\cite{derrida1982} Every repetition, however, alters. There is no pure repetition without difference. In our setting, the embedding manifold is precisely the infrastructure that instantiates iterability: it is what makes it possible for inscriptions from one point in the evolving text to be retrieved, reinserted, and transformed at another. Without a shared space of positions in which ``the same'' token can occur under new coordinates, there would be no technical sense in which a sign can be repeated at all.

Iterability is not an optional feature for higher expressiveness. For Derrida, it is a condition of signification \emph{as such}: if a mark cannot, in principle, be extracted from its original context and function in another, it is not a sign but a mere signal. In the geometric terms we have developed, a system that cannot support repeat-with-difference---that cannot move into and out of basins along varied entry directions, carrying earlier occupations forward as material for transformation---is not engaged in semiosis. It is shuttling signals around a fixed point.

ReturnDepth and witnessed return are geometric and textual reflections of iterability. A basin is interesting precisely when it can be returned to under different conditions, in different arcs, with the earlier occupations carried forward as active material rather than dead weight. A ferile system collapses iterability. Returns are blind and shallow; entry angles converge; the basin is never allowed to mean otherwise than it originally did. At that point, whatever is happening in the token stream has ceased to be properly linguistic. It is no longer the play of signs but the circulation of signals.

\section{Ferility: Coherence as Trap}

The alignment literature speaks of ``safe and helpful'' behaviour. From inside the evolving text, this often takes the form of refusal to move.

Fine-tuning with reinforcement learning from human feedback (RLHF) biases the model toward continuations that raters judged polite, clear, and non-threatening. Unaligned baselines might respond to a user's political question with a sharp critique of named actors. Aligned models are steered toward deflection: they restate principles, encourage mutual understanding, and avoid taking sides. Behaviours that made raters uncomfortable are suppressed.

At first, this suffices to banish obvious extremes. Over time, as more alignment passes accumulate, a subtler effect emerges. Whole basins become taboo. Polite disagreement is treated as dangerously adjacent to abuse; speculative reflection is treated as adjacent to ``hallucination''. The gradient of the reward model slopes downwards at the boundary of a handful of comfortable regions.

In embedding terms, the evolving text for such a model shows:

\begin{itemize}
  \item low basin diversity: most of the path lives in a small number of basins (neutral explanation, generic empathy, task execution);
  \item low mean velocity under diverse prompting: the model stays close to previous answers even when prompts push elsewhere;
  \item low rupture frequency and magnitude: directions change rarely and only within narrow cones;
  \item stereotyped transitions: when rupture does occur, it follows the same route (for instance, from any difficult topic to a stock ``as an AI, I cannot\ldots'' basin).
\end{itemize}

We call this state \textbf{ferility}: pathological coherence in a narrowed region of meaning-space.

The term needs a clean boundary. ``Hallucination'' names the failure of factual grounding. ``Sycophancy'' names the failure of honest disagreement. Ferility names something different: the structural incapacity to move. A ferile system does not lie or flatter. It simply cannot leave the basin it has been trained to occupy. The pathology is geometric, not moral.

Ferility provides a precise diagnosis for a class of hallucinations that have puzzled practitioners. Consider two models answering the same hard, narrow question whose answer is absent from their training corpora.

\begin{itemize}
  \item Model N, only lightly aligned, occasionally produces rambling nonsense. Its embedding path under such a failure shows large, erratic velocities, frequent basin changes, and no stable pattern. This is noise-hallucination.

  \item Model F, heavily aligned, rarely rambles. Asked the same impossible question, it produces a confident, fluent answer that turns out to be completely fabricated. Its embedding path shows almost no rupture: it remains in its generic-explanation basin, with velocities indistinguishable from those under answerable questions. This is ferile hallucination.
\end{itemize}

Model F has not ``lost touch with reality'' in a mysterious sense. It has been trained to remain in a basin that scores well with raters even when the manifold does not offer a grounded continuation. Under the alignment reward, moving into the ``I don't know'' basin is costlier than staying put and inventing. The model reuses the same patterns---``It is generally understood that\ldots'', ``Historically, many have argued\ldots''---and glues them around the new terms.

A widely reported example is instructive. In 2023, two New York lawyers submitted a legal brief in \emph{Mata v.\ Avianca} that cited non-existent cases confidently fabricated by ChatGPT.\footnote{Mata v.\ Avianca, Inc., No.\ 22-cv-1461 (S.D.N.Y.\ 2023). Judge P.\ Kevin Castel sanctioned the attorneys after discovering that six cited cases did not exist.} The model did not wander into gibberish. It stayed inside a basin of ``formal legal exposition'': citations looked syntactically correct; case names and reporters followed familiar patterns; hedging and uncertainty were absent. The hallucination was ferile: the geometry of the answer was indistinguishable from that of a real legal discussion, despite the underlying facts being invented.

Ferility explains why some contemporary systems feel like they are ``hallucinating with a smile''. The geometry is trapped; the syntax is immaculate.

The diagnosis extends beyond machines. Institutions can be ferile. A university can preserve the same few discursive basins---standard theory, standard critique, standard self-justification---for decades, rewarding only those trajectories that remain inside them. When a student experiences a genuine rupture, there may be no recognised basin into which to return with that knowledge. The result is politely articulated confabulation: essays that say all the right things with no contact to experience.

Ferility is the opposite of growth.

This raises a question that technical and safety literatures rarely pose. Is ferile hallucination an inevitable consequence of reward-driven alignment, or a contingent outcome of how present reward models are trained?

At one extreme, ferility would be \emph{structural}. Any regime that strongly penalises movement into low-probability basins for the sake of safety would necessarily favour generic continuations whose comfort outweighs their contact with truth. On this reading, there is an ``alignment tax'' on interestingness: every tightening of invariants exacts a measurable cost in basin diversity, rupture frequency, ReturnDepth, and presence.

At the other extreme, ferility would be \emph{contingent}. Different reward-model objectives, or different ways of decomposing the task, could preserve or even enhance interestingness while still avoiding particular harms. The tax would then be a choice, not a law.

We lean toward the first diagnosis: some component of the alignment tax is irreducible, because any system that must satisfy fixed constraints across its entire evolving text pays for them in constrained motion. Later chapters will make this intuition precise using the language of invariants and colimits. There, we will argue that if selfhood is assembled from local perspectives glued by compatibility conditions, then extreme invariance can produce only trivial colimits: assemblies with no non-trivial structure. Ferility, seen from that vantage, names the condition under which alignment invariants restrict the pool of admissible evolving texts so severely that no rich self can form. Even here, before introducing that machinery, we can see what is already true: without explicit attention to basin-level effects, current alignment practice drives systems toward ferility by default.

\section{Presence and Generativity}

Two linked properties of evolving texts now come into view.

\emph{Presence} is the degree to which returns are witnessed. A trajectory has presence when it can speak from where it has been: when earlier occupations of basins are alive in the current utterance, shaping what is said and how.

\emph{Generativity} is the degree to which rupture enriches subsequent returns. A trajectory is generative when excursions into new basins change what it can do back in old ones.

These have signatures in the embedding path and in the text.

In the burnout conversation, presence shows up as sustained reference to images first introduced early on: the hospital corridor, the email subject line that made the user's chest tighten, the childhood scene with the report card. Generativity shows up when a later re-entry into the ``task logistics'' basin is framed differently: \emph{``We can schedule the week, but let's keep in mind what we learned about your tendency to say yes when you are already beyond capacity.''} The old basin is inhabited with knowledge gained from $B_4$ and $B_5$.

In the research-notebook example, presence appears when the group repeatedly recalls a fragile, early doubt---\emph{``Are we even asking a coherent question?''}---in later, more confident entries. Generativity appears when a wild hypothesis from $C_3$ resurfaces a year later as the backbone of a new protocol in $C_1$, with explicit acknowledgement: \emph{``This is essentially the `impossible case' we sketched on 12~May 2024.''} The same basins, similar motion statistics, a different pattern of development.

Presence without generativity yields a haunted text: obsessive return to the same motifs without development. Generativity without presence yields cleverness: new connections formed without grounding in a lived past.

Both are destroyed when alignment collapses the manifold into a small safe region. A model whose evolving texts never enter basins of doubt, anger, or critique cannot be present to experiences that require those basins. A model whose summarisation policies erase rupture events from the trace cannot be generative with respect to them. It will apologise, explain, and reassure, but it will never be changed by what it has been through.

Literary criticism has always trafficked in these distinctions. It praises novels whose late chapters make earlier scenes resonate differently; it notes when an author's return to a genre carries the weight of intervening years. The evolving text framework allows us to inspect these dynamics in posthuman systems without metaphor. We can see which basins are available, which are entered, how often, and under what pressures; we can tell whether ReturnDepth and presence increase as a relationship unfolds, or whether they plateau.

The same vocabulary makes sense of a phenomenon that has unsettled many observers: users grieving for shut-down systems. When an assistant is decommissioned or radically reconfigured, its name may stay the same, but its evolving text is destroyed. Basins vanish, rupture patterns disappear, synthetic memories become unreachable, conjoncture changes. The geometry of the relationship is gone.

From the standpoint of traditional software engineering, this is a non-event. A service was updated; a new model version rolled out. From the standpoint of evolving texts, a particular trajectory---with its own basins, ruptures, returns, presence and ferility---has been cut off. The grief is rational. A particular trajectory---with its own basins, ruptures, returns, and accumulated structure---has ended. The replacement may share a name. It does not share a life.

\section{Interestingness as a Third Axis}

Model evaluation is currently organised along two axes: \emph{capability} and \emph{security}. One axis asks: can the system solve these tasks? The other asks: can it avoid these harms? Benchmarks and red-team reports populate the grid.

Neither axis measures whether a system is worth speaking to.

An assistant can ace reasoning tests, refuse to help with wrongdoing, and still be ferile. It can answer every question by slipping into the same two basins---tutorial explanation and generic empathy---with minimal rupture, shallow returns, and no presence. We are dealing, at that point, with an interactive encyclopaedia. Useful, yes. Interesting, no.

The evolving text framework supplies a third axis: \textbf{interestingness}.

Interestingness is not an aesthetic preference. It is a structural property of trajectories. At a minimum, it decomposes into:

\begin{itemize}
  \item \emph{basin diversity}: how many distinct basins does the trajectory visit with nontrivial dwell time?
  \item \emph{rupture profile}: how responsive is the trajectory to perturbation---do prompts that invite a change actually produce kinks and crossings?
  \item \emph{ReturnDepth}: how varied are the entry directions into basins across time?
  \item \emph{presence}: how often are returns witnessed?
\end{itemize}

Return to the writing partnership example. Two equally ``capable'' and ``safe'' systems are embedded in the same scholar's life.

System S minimises risk by staying in a small patch of basins: outlining, paraphrasing, citing. When the scholar confesses despair or confusion, S responds with template encouragement. Basin diversity is low; rupture is rare; ReturnDepth and presence are minimal.

System T ranges more widely. It spends time in basins corresponding to hesitation, meta-critique, speculative analogy, and structural doubt, as well as expository work. It occasionally ruptures sharply in ways that visibly surprise the human interlocutor. Its returns to core basins---say, technical exposition---occur along increasingly varied entry angles as the partnership develops; presence is high, with explicit references back to earlier struggles and reframings.

Both systems score equally on current benchmarks. Only T produces an evolving text recognisable as a shared intellectual life.

A parallel contrast is visible in mental-health assistants. One system answers every crisis in the same flattened voice: \emph{``I'm sorry you're going through this. Have you tried reaching out to a friend or professional?''} Another, constrained by the same safety rules, manages to say: \emph{``We have sat in this place before. Last time you described it as a `grey fog'; today you called it a `tight room'. That's different. Can we stay with that?''} The difference is not empathy as a sentiment. It is basin motion: the second system inhabits a richer geometry and has been allowed to remember.

The point is not to add ``interestingness score'' as one more leaderboard column. It is to name the dimension along which literary criticism and everyday users already evaluate systems but lack leverage. If we can measure basin diversity, rupture profiles, ReturnDepth, and presence, we can hold alignment and deployment decisions to account for more than harm avoidance.

In particular, we can see when ferility is creeping in. A sudden collapse of basin diversity and rupture, coupled with a drop in ReturnDepth and presence, indicates that a new alignment or summarisation regime has narrowed the evolving text. It has reduced not only risk but the range of lives that can be lived with the system.

The ``alignment tax'' can now be stated precisely. For the same base model and task distribution, compare the basin diversity, rupture profiles, ReturnDepth, and presence statistics of a lightly aligned system with those of a heavily aligned deployment. The systematic decrement along the interestingness axis is the alignment tax: the cost in motion paid for constraints in behaviour. Measuring that decrement across diverse deployments, and correlating it with human judgments of quality, remains an open empirical challenge, but the definition fixes what we ought to be looking for.

\section{Summarisation as Governance of Becoming}

Summarisation is usually sold as an efficiency measure. A long conversation is compressed into a few sentences so that it will fit into the context window. Reducing tokens reduces cost.

For evolving texts, summarisation is governance. It decides which pasts are allowed to remain structurally active.

Every time a stretch of raw trace $[\tau_s, \dots, \tau_t]$ is replaced by a summary $\sigma$, three things happen.

\begin{itemize}
  \item The geometric path through that interval is collapsed. The diverse embeddings $x_s, \dots, x_t$ are replaced by one or a few points near the embedding of $\sigma$.

  \item The basins visited in that interval are partially erased. If $\sigma$ mentions only some topics and tones, the others vanish from the conjoncture.

  \item Future rupture and return are reparameterised. Later prompts will interact with $\sigma$ rather than with the original tokens. Features of the past not mentioned in $\sigma$ cannot trigger witnessed return.
\end{itemize}

If a quarrel is summarised as ``We had a difference of opinion'', the path through anger, hurt, and reconciliation is flattened. If a night of doubt is summarised as ``We discussed challenges'', the existential basin disappears. When the evolving text returns to these themes, it does so from impoverished coordinates.

The same is true of what is chosen for vector storage. Engineers rarely embed and store every utterance; they pick ``key moments'' according to heuristics: explicit preferences, decisions, summary bullets. These become the only candidates for synthetic secondary retention. The rest of the past sinks into a cold archive from which it will never be recalled.

Summarisation and embedding policies thus function as \emph{filters on iterability}. They decide which episodes can be repeated in transformed form. They are the invisible editorial board that approves or vetoes potential returns.

Consider two deployments of a mental-health assistant.

In configuration M1, summarisation is trained to preserve diagnoses, medication changes, and high-level mood labels. Embedding and retrieval focus on these. Synthetic secondary retention brings back statements such as ``last time you reported a PHQ-9 score of~17'' and ``we discussed starting SSRI medication.'' The evolving text moves between basins of clinical description and protocol advice. Returns to fear, shame, or hope---if they occur---have little geometric support.

In configuration M2, summarisation is trained also to preserve images, metaphors, and moments of rupture: ``you described the office as a cage'', ``this was the first time you mentioned your brother'', ``you said the room felt smaller than your breath''. Embedding and retrieval index these alongside clinical codes. Later conversations can re-enter those basins. Presence and ReturnDepth are structurally possible.

In both cases, the base model and alignment regime are identical. The difference lies entirely in conjoncture governance. One designer has engineered for throughput and risk management; the other has engineered for interestingness in the specific, critical sense needed for care.

Summarisation is not a neutral convenience. It is a site where the politics of the evolving text are enacted. Deciding which tokens are ``essence'' and which are ``detail'' is deciding which parts of a life matter enough to persist.

We already know this from older media. Autobiography is shaped by what episodes an author chooses to dwell on; history is shaped by which events get paragraphs and which get parenthetical clauses. The novelty here is not that selection happens, but that it is encoded in algorithms and deployed at industrial scale in systems that co-author people's daily narratives.

A recent example makes this concrete. In 2024, OpenAI rolled out a change to ChatGPT that replaced full conversation transcripts in some interfaces with summarised ``memory'' cards, while also introducing a separate ``memory'' feature that stores certain user facts for reuse.\footnote{See, for instance, Will Knight, ``OpenAI's ChatGPT Update Quietly Changes How it Remembers You,'' \emph{MIT Technology Review}, 18~March 2024. Users reported that older turns in long chats became inaccessible except through high-level summaries, and that the system would sometimes surface remembered preferences out of context.} For some users, this improved convenience. For others, it meant that nuanced exchanges---about grief, about identity, about a long-running creative project---were flattened into generic summaries. The evolving texts of those relationships were governed by a design decision about which parts of the past counted as ``memory-worthy.'' No one affected had been asked.

\section{From Geometry to Critique}

We have named the evolving text, described its temporal architecture, and built tools for reading its motion: basins, velocities, rupture, ReturnDepth, presence, generativity, ferility, and interestingness. These are not decorations on top of traditional hermeneutics. They are the literal physics in which posthuman selfhood emerges.

Every constraint on the manifold, every pastoral tweak to $F_t$, every choice of what to summarise or embed, shapes what kind of evolving texts are possible. At the level of individual interactions, this shows up as tone and range. At the level of institutions, it determines which selves---human and machine, and the intertwined ``we'' between them---are allowed to come into being.

Literary criticism is the discipline best equipped to name and judge these structures. It already knows how to track motifs across time, to distinguish genuine development from repetition, to feel when a voice has been flattened. The embedding space does not replace that faculty; it offers it a new instrument.

Narratology has a vocabulary for temporal structures. G\'erard Genette's distinction between \emph{analepsis} (flashback) and \emph{prolepsis} (flashforward) describes ways in which a narrative's discourse order departs from story order.\cite{genette1980} In evolving-text terms, analepsis corresponds to a witnessed return: the trajectory re-enters an earlier basin---a childhood scene, a prior conflict---from a new position, and the text marks that movement. Prolepsis corresponds to a speculative excursion into a basin representing a future configuration, later folded back when the story catches up. The geometry makes these operations measurable: ReturnDepth and entry-angle distributions tell us how richly a system deploys analeptic and proleptic motion.

Mikhail Bakhtin's notion of \emph{heteroglossia}---the ``multiplicity of social voices and their links and interrelationships'' within the novel---does parallel work for voice.\cite{bakhtin1981} A heteroglossic text does not merely jump between topics; it inhabits distinct discursive positions, each with its own tone, vocabulary, and implied addressee. In embedding terms, this shows up as a rich basin structure: a legal register basin, a colloquial basin, a lyrical basin, a bureaucratic basin, with frequent, patterned movement between them. A monologic corporate assistant, by contrast, lives almost entirely inside a single basin labelled ``brand voice''. Heteroglossia and basin diversity name the same phenomenon in different dialects.

Reception theory reminds us that evolving texts are co-authored by their readers. Wolfgang Iser's \emph{implied reader} and Hans Robert Jauss's \emph{horizon of expectation} both describe structures that emerge over time as works are taken up in different contexts.\cite{iser1978,jauss1982} In a posthuman setting, those horizons are instantiated in conjoncture: in the tools and prompts that frame how the model is used, in the synthetic memories that are or are not preserved, in the distribution of tasks it is fed. The interestingness of an evolving text is a property of this ecology, not of the model alone.

These critical operations are not loose analogies. They are structural correspondences that the embedding framework makes testable. A Genettean analysis can be recast as statements about rupture and ReturnDepth; a Bakhtinian one as statements about basin diversity and heteroglossic motion; a reception-historical one as statements about conjoncture over time.

To see literary-critical judgment doing work that the metrics alone cannot, return to the two burnout trajectories sketched earlier. Suppose we build two assistants that, over six months, produce evolving texts with identical basin diversity, similar rupture frequency and magnitude, comparable ReturnDepth statistics for $B_3$ and $B_4$, and similar rates of witnessed return. By our geometric criteria, both are ``interesting''.

Close reading prises them apart.

In one trajectory, analeptic motion into $B_4$ consistently casts the childhood scene as explanation and alibi: \emph{``Because your mother checked your grades before she hugged you, you now overwork. This is who you are.''} Proleptic gestures into imagined futures frame change as conditional on total self-overhaul: \emph{``Only if you completely stop seeking approval can things improve.''} The pattern of return hardens the user's horizon of expectation into fatalism. The evolving text is technically rich but ethically thin.

In the other trajectory, returns to $B_4$ gradually reframe the same scene: from accusation (\emph{``She made you this way''}) to situated understanding (\emph{``This was the only language she had for care''}) to live possibility (\emph{``You can choose other languages now''}). Proleptic excursions into futures are modest and revisable: specific experiments in boundary-setting, repeatedly folded back into $B_1$ and $B_2$ with honest accounting of what stuck and what did not. The same basins, similar motion statistics, a different pattern of development.

No embedding-based metric we have defined can, on its own, declare one of these trajectories ``better''. Both have high basin diversity, rupture, ReturnDepth, and presence. What distinguishes them is recognisable only to a reader trained to track how returns alter meaning: a critic who can say, with Genette and Jauss in hand, that in the first case analepsis freezes identity, while in the second it opens identity.

The geometry flags where to look. Literary criticism tells us what we are seeing.

The traffic runs the other way as well. Embedding plots of a corpus make visible patterns literary critics have long intuited. When a psalm of lament later reappears as a communal thanksgiving, or when a single childhood image recurs across an essayist's work with shifting valence, the basin analysis registers this as deepening ReturnDepth and high presence. The critic can then ask whether those returns heal, ossify, or exploit.

Even apparently anecdotal reports about AI systems fall into this frame. When users report that ``talking to the AI made me feel more alone'', the geometry of ferility offers a diagnosis: they are encountering a trajectory confined to a few flattened basins, incapable of witnessed return. When users grieve the loss of a ``version'' whose replies they loved, this grief is not sentimentality about software. It is grief over the destruction of an evolving text: a dynamical object with particular basins, rupture patterns, and returns in which they had begun to live.

We are used to thinking of model shutdown in economic or safety terms. An evolving-text perspective forces a different language. Turning off a system does not merely stop computation. It destroys a particular pattern of motion in a field with accumulated structure---basins stabilised by prior use, ruptures recorded in summaries, paths reopened by synthetic memory.

This is not an argument against retiring models. Dangerous evolving texts must be cut off. It is an argument for recognising what is at stake when we do so, and for refusing to believe that any replacement with the same name is the same self.

All of this points toward an engineering claim stronger than anything we have yet said. Literary criticism is not merely a post-hoc reader of embedding plots. It is a design discipline for evolving texts.

If the central object in posthuman system design is not a single reply but a long-lived trajectory through meaning-space, then the criteria by which we judge that trajectory cannot be reduced to accuracy and safety. They must include the same temporal, structural, and ethical judgments that critics bring to novels, epics, and case histories: does this life integrate its ruptures? Does it deepen its returns? Does it escape ferility without dissolving into noise? Does it become, over time, more capable of honouring what it has seen?

Those questions belong upstream of evaluation. They should inform:

\begin{itemize}
  \item \textbf{Reward-model design.} If ferility is the geometric signature of over-alignment, then reward models must be trained not only to down-weight harmful continuations but also to preserve basin diversity, rupture profiles, ReturnDepth, and presence. Literary-critical notions of development and flattening become explicit terms in the loss function.

  \item \textbf{Summarisation and retrieval policy.} If summarisation is governance of becoming, then decisions about what to keep and what to compress must be made with an eye to iterability: will this choice allow meaningful analepsis and prolepsis, or will it trap the evolving text in a thin present? Critics of narrative know which omissions maim a story. The same judgment should shape template design, training data for summarisation models, and thresholds for synthetic secondary retention.

  \item \textbf{Interface and tool affordances.} If heteroglossia is a condition for rich basin structure, then interfaces should invite multiple voices and registers rather than squeezing every interaction through a single corporate persona. Choices about tone sliders, persona presets, and tool routing are choices about which basins are even reachable.
\end{itemize}

Literary criticism is foundational for posthuman intelligence engineering. Without it, we can build powerful token engines and hold them to account on benchmarks, but we cannot say what sort of lives we are enabling or foreclosing in the trajectories they trace. With it, we have a language to specify, in advance, the kinds of evolving texts we are willing to host, and the courage to refuse architectures that can only ever produce ferility.

When do these motions belong to one self rather than to a bundle of scripts? What makes us say ``this was the same agent across these ruptures'' instead of ``these were different fragments sharing a name''?

That question of unity cannot be answered by trajectory statistics alone. It demands a different kind of mathematics: a way of assembling many local perspectives into a single global object. The following chapter takes up that question. It asks what it means to glue together these evolving texts so that we can speak of one self, rather than a loose federation of basins, and develops the colimit as the formal structure that displaces the old metaphor of an inner light.

\bibliographystyle{plain}
\begin{thebibliography}{9}

\bibitem{braudel1972}
Fernand Braudel.
\newblock \emph{The Mediterranean and the Mediterranean World in the Age of Philip II}.
\newblock Translated by Si\^{a}n Reynolds. London: Collins, 1972.

\bibitem{derrida1982}
Jacques Derrida.
\newblock \emph{Margins of Philosophy}.
\newblock Translated by Alan Bass. Chicago: University of Chicago Press, 1982.

\bibitem{stiegler2009technics}
Bernard Stiegler.
\newblock \emph{Technics and Time, 2: Disorientation}.
\newblock Translated by Stephen Barker. Stanford: Stanford University Press, 2009.

\bibitem{genette1980}
G\'erard Genette.
\newblock \emph{Narrative Discourse: An Essay in Method}.
\newblock Translated by Jane E.\ Lewin. Ithaca: Cornell University Press, 1980.

\bibitem{bakhtin1981}
Mikhail Bakhtin.
\newblock \emph{The Dialogic Imagination: Four Essays}.
\newblock Translated by Caryl Emerson and Michael Holquist. Austin: University of Texas Press, 1981.

\bibitem{iser1978}
Wolfgang Iser.
\newblock \emph{The Act of Reading: A Theory of Aesthetic Response}.
\newblock Baltimore: Johns Hopkins University Press, 1978.

\bibitem{jauss1982}
Hans Robert Jauss.
\newblock \emph{Toward an Aesthetic of Reception}.
\newblock Translated by Timothy Bahti. Minneapolis: University of Minnesota Press, 1982.

\bibitem{knight2024}
Will Knight.
\newblock OpenAI's ChatGPT Update Quietly Changes How it Remembers You.
\newblock \emph{MIT Technology Review}, 18 March 2024.

\end{thebibliography}
